\chapter{2009年北京卷（理）}

\section{单选题}

\begin{problem}
在复平面内，复数 \(z = \i(1 + 2\i)\) 对应的点位于（\quad）

\choices
    {第一象限}
    {第二象限}
    {第三象限}
    {第四象限}

\begin{answer}
B
\end{answer}

\begin{solution}
将复数化为代数形式，得 \(z=\i(1+2\i)=\i+2\i^2=-2+\i\). 因而它在复平面内对应的点为 \((-2,1)\)，实部为负、虚部为正，位于第二象限，故选 B.
\end{solution}
\end{problem}

\begin{problem}
已知向量 \(\bs{a}\)，\(\bs{b}\) 不共线，\(\bs{c} = k\bs{a} + \bs{b}\)（\(k \in \R\)），\(\bs{d} = \bs{a} - \bs{b}\). 如果 \(\bs{c} \parallel \bs{d}\)，那么（\quad）

\choices
    {\(k = 1\) 且 \(\bs{c}\) 与 \(\bs{d}\) 同向}
    {\(k = 1\) 且 \(\bs{c}\) 与 \(\bs{d}\) 反向}
    {\(k = -1\) 且 \(\bs{c}\) 与 \(\bs{d}\) 同向}
    {\(k = -1\) 且 \(\bs{c}\) 与 \(\bs{d}\) 反向}

\begin{answer}
D
\end{answer}

\begin{solution}
因为 \(\bs{a}\)，\(\bs{b}\) 不共线，而 \(\bs{c}\parallel\bs{d}\)，所以存在实数 \(\lambda\) 使 \(k\bs{a}+\bs{b}=\lambda(\bs{a}-\bs{b})\). 比较 \(\bs{a}\) 和 \(\bs{b}\) 的系数，得 \(k=\lambda\)，\(1=-\lambda\)，从而 \(k=-1\). 此时 \(\bs{c}=-\bs{a}+\bs{b}=-(\bs{a}-\bs{b})=-\bs{d}\)，所以两向量反向，故选 D.
\end{solution}
\end{problem}

\begin{problem}
为了得到函数 \( y = \lg \frac{x + 3}{10} \) 的图象，只需把函数 \( y = \lg x \) 的图象上所有的点（\quad）

\choices
    {向左平移 \(3\) 个单位长度，再向上平移 \(1\) 个单位长度}
    {向右平移 \(3\) 个单位长度，再向上平移 \(1\) 个单位长度}
    {向左平移 \(3\) 个单位长度，再向下平移 \(1\) 个单位长度}
    {向右平移 \(3\) 个单位长度，再向下平移 \(1\) 个单位长度}

\begin{answer}
C
\end{answer}

\begin{solution}
由对数运算，\(\lg\dfrac{x+3}{10}=\lg(x+3)-1\). 将 \(y=\lg x\) 的图象向左平移 \(3\) 个单位长度得到 \(y=\lg(x+3)\)，再向下平移 \(1\) 个单位长度得到所求图象，故选 C.
\end{solution}
\end{problem}

\begin{problem}
若正四棱柱 \(ABCD - A_{1}B_{1}C_{1}D_{1}\) 的底面边长为 \(1\)，\(AB_{1}\) 与底面 \(ABCD\) 成 \(60^{\circ}\) 角，则 \(A_{1}C_{1}\) 到底面 \(ABCD\) 的距离为（\quad）

\choices
    {\(\frac{\sqrt{3}}{3}\)}
    {\(1\)}
    {\(\sqrt{2}\)}
    {\(\sqrt{3}\)}

\begin{answer}
D
\end{answer}

\begin{solution}
设正四棱柱的高为 \(h=BB_1\). 由于 \(AB_1\) 在底面上的射影为 \(AB\)，且 \(AB=1\)，所以 \(\tan 60^\circ=\dfrac{BB_1}{AB}=h\)，得 \(h=\sqrt3\). 直线 \(A_1C_1\) 平行于底面 \(ABCD\)，且位于距底面 \(h\) 的平面内，因此 \(A_1C_1\) 到底面 \(ABCD\) 的距离为 \(\sqrt3\)，故选 D.
\end{solution}
\end{problem}

\begin{problem}
“ \( \alpha=\frac{\pi}{6}+2k\pi \)”（\(k\in\Z\)）是“ \( \cos2\alpha=\frac{1}{2} \) ”的（\quad）

\choices
    {充分而不必要条件}
    {必要而不充分条件}
    {充分必要条件}
    {既不充分也不必要条件}

\begin{answer}
A
\end{answer}

\begin{solution}
若 \(\alpha=\dfrac\pi6+2k\pi\)，则 \(2\alpha=\dfrac\pi3+4k\pi\)，从而 \(\cos2\alpha=\cos\dfrac\pi3=\dfrac12\)，所以前一个条件是后一个条件的充分条件.

反之，若 \(\cos2\alpha=\dfrac12\)，则 \(2\alpha=2n\pi\pm\dfrac\pi3\)，即 \(\alpha=n\pi\pm\dfrac\pi6\)（\(n\in\Z\)）. 这包含不满足 \(\alpha=\dfrac\pi6+2k\pi\) 的情形，例如 \(\alpha=\dfrac{5\pi}{6}\)，故前一个条件不是必要条件. 因此它是充分而不必要条件，故选 A.
\end{solution}
\end{problem}

\begin{problem}
若 \(\left(1 + \sqrt{2}\right)^5 = a + b\sqrt{2}\) （\(a\)，\(b\) 为有理数），则 \(a + b =\)（\quad）

\choices
    {\(45\)}
    {\(55\)}
    {\(70\)}
    {\(80\)}

\begin{answer}
C
\end{answer}

\begin{solution}
利用二项式展开并合并有理项与含 \(\sqrt2\) 的项：\((1+\sqrt2)^5=1+5\sqrt2+10\cdot2+10\cdot2\sqrt2+5\cdot4+4\sqrt2=41+29\sqrt2\). 因而 \(a=41\)，\(b=29\)，所以 \(a+b=70\)，故选 C.
\end{solution}
\end{problem}

\begin{problem}
用 \(0\) 到 \(9\) 这 \(10\) 个数字，可以组成没有重复数字的三位偶数的个数为（\quad）

\choices
    {\(324\)}
    {\(328\)}
    {\(360\)}
    {\(648\)}

\begin{answer}
B
\end{answer}

\begin{solution}
按个位数字分类. 当个位为 \(0\) 时，百位有 \(9\) 种选法，十位从剩余数字中有 \(8\) 种选法，共有 \(9\times8=72\) 个.

当个位为 \(2\)，\(4\)，\(6\)，\(8\) 时，个位有 \(4\) 种选法；百位不能为 \(0\)，且不能与个位重复，有 \(8\) 种选法；十位从剩余数字中有 \(8\) 种选法，共有 \(4\times8\times8=256\) 个. 因此总数为 \(72+256=328\)，故选 B.
\end{solution}
\end{problem}

\begin{problem}
点 \(P\) 在直线 \(l: y = x - 1\) 上，若存在过 \(P\) 的直线交抛物线 \(y = x^2\) 于 \(A\)，\(B\) 两点，且 \(|PA| = |AB|\)，则称点 \(P\) 为“\(A\) 点”，那么下列结论中正确的是（\quad）

\choices
    {直线 \(l\) 上的所有点都是“\(A\) 点”}
    {直线 \(l\) 上仅有有限个点是“\(A\) 点”}
    {直线 \(l\) 上的所有点都不是“\(A\) 点”}
    {直线 \(l\) 上有无穷多个点（但不是所有的点）是“\(A\) 点”}

\begin{answer}
A
\end{answer}

\begin{solution}
任取直线 \(l\) 上的点 \(P=(p,p-1)\). 在抛物线 \(y=x^2\) 上取 \(A=(m,m^2)\)，并令 \(B=2A-P\)，则 \(A\) 是线段 \(PB\) 的中点，从而 \(P\)，\(A\)，\(B\) 共线且 \(|PA|=|AB|\).

只需使 \(B\) 在抛物线上. 由 \(B=(2m-p,2m^2-p+1)\)，条件为 \(2m^2-p+1=(2m-p)^2\)，即 \(2m^2-4pm+p^2+p-1=0\). 这是关于 \(m\) 的二次方程，其判别式为 \(\Delta=(-4p)^2-8(p^2+p-1)=8(p^2-p+1)>0\)，所以对任意实数 \(p\) 都有实数解 \(m\). 此时 \(B\) 在抛物线上. 若 \(A=B\)，则 \(P=A\)，从而 \(p^2-p+1=0\)，但该方程的判别式为 \(-3<0\)，矛盾，故 \(A\ne B\). 因此每个 \(P\in l\) 都是“\(A\) 点”，故选 A.
\end{solution}
\end{problem}

\section{填空题}

\begin{problem}
\(\displaystyle\lim_{x\to 1}\frac{x\sqrt{x} - x}{x - 1} =\)\fillinblank{}.

\begin{answer}
\(\frac{1}{2}\)
\end{answer}

\begin{solution}
当 \(x\to1\) 时，\(x\geqslant0\)，且 \(\dfrac{x\sqrt{x}-x}{x-1}=\dfrac{x(\sqrt{x}-1)}{(\sqrt{x}-1)(\sqrt{x}+1)}=\dfrac{x}{\sqrt{x}+1}\). 因而原极限为 \(\dfrac{1}{1+1}=\dfrac12\).
\end{solution}
\end{problem}

\begin{problem}
若实数 \(x\)，\(y\) 满足 \(\begin{cases}
x + y - 2 \geqslant 0,\\
x \leqslant 4,\\
y \leqslant 5.
\end{cases}\) 则 \(s = y - x\) 的最小值为\fillinblank{}.

\begin{answer}
\(-6\)
\end{answer}

\begin{solution}
由 \(x+y-2\geqslant0\) 得 \(y\geqslant2-x\)，所以 \(s=y-x\geqslant2-2x\). 又因 \(x\leqslant4\)，有 \(2-2x\geqslant2-8=-6\)，故 \(s\geqslant-6\).

取 \(x=4\)，\(y=-2\)，则三个条件均满足，且 \(s=y-x=-2-4=-6\). 因此 \(s\) 的最小值为 \(-6\).
\end{solution}
\end{problem}

\begin{problem}
设 \( f(x) \) 是偶函数. 若曲线 \( y = f(x) \) 在点 \( (1, f(1)) \) 处的切线的斜率为 \(1\)，则该曲线在 \((-1, f(-1))\) 处的切线的斜率为\fillinblank{}.

\begin{answer}
\(-1\)
\end{answer}

\begin{solution}
偶函数满足 \(f(-x)=f(x)\). 对等式两边求导，得 \(-f'(-x)=f'(x)\). 令 \(x=1\)，则 \(-f'(-1)=f'(1)=1\)，所以 \(f'(-1)=-1\)，即所求切线斜率为 \(-1\).
\end{solution}
\end{problem}

\begin{problem}
椭圆 \(\frac{x^2}{9} + \frac{y^2}{2} = 1\) 的焦点为 \(F_1\)，\(F_2\)，点 \(P\) 在椭圆上. 若 \(|PF_1| = 4\)，则 \(|PF_2| = \)\fillinblank{}；\(\angle F_1PF_2\) 的大小为\fillinblank{}.

\begin{answer}
\(2\)，\(120^{\circ}\)
\end{answer}

\begin{solution}
椭圆的长半轴为 \(a=3\)，短半轴平方为 \(b^2=2\)，故半焦距满足 \(c^2=a^2-b^2=7\)，从而 \(|F_1F_2|=2\sqrt7\). 由椭圆定义，\(|PF_1|+|PF_2|=2a=6\)，所以 \(|PF_2|=6-4=2\).

在 \(\triangle F_1PF_2\) 中，由余弦定理，\(\cos\angle F_1PF_2=\dfrac{4^2+2^2-(2\sqrt7)^2}{2\cdot4\cdot2}=-\dfrac12\). 因此 \(\angle F_1PF_2=120^\circ\).
\end{solution}
\end{problem}

\begin{problem}
若函数 \( f(x) = \begin{cases}
\frac{1}{x}, & x < 0,\\
\left(\frac{1}{3}\right)^x, & x \geqslant 0,
\end{cases} \) 则不等式 \( |f(x)| \geqslant \frac{1}{3} \) 的解集为\fillinblank{}.

\begin{answer}
\([-3,1]\)
\end{answer}

\begin{solution}
当 \(x<0\) 时，\(f(x)=\dfrac1x<0\)，故 \(|f(x)|=-\dfrac1x\). 不等式 \(-\dfrac1x\geqslant\dfrac13\) 在 \(x<0\) 下等价于 \(x\geqslant-3\)，所得解集为 \([-3,0)\).

当 \(x\geqslant0\) 时，\(|f(x)|=\left(\dfrac13\right)^x\). 因为底数 \(\dfrac13\in(0,1)\)，所以 \(\left(\dfrac13\right)^x\geqslant\left(\dfrac13\right)^1\) 等价于 \(x\leqslant1\)，所得解集为 \([0,1]\). 合并得解集 \([-3,1]\).
\end{solution}
\end{problem}

\begin{problem}
已知数列 \(\{a_{n}\}\) 满足：\(a_{4n - 3} = 1\)，\(a_{4n - 1} = 0\)，\(a_{2n} = a_{n}\)，\(n \in \N^*\)，则 \(a_{2009} = \)\fillinblank{}；\(a_{2014} = \)\fillinblank{}.

\begin{answer}
\(1\)，\(0\)
\end{answer}

\begin{solution}
由 \(2009=4\times503-3\)，代入 \(a_{4n-3}=1\) 得 \(a_{2009}=1\). 又由 \(a_{2n}=a_n\)，有 \(a_{2014}=a_{1007}\). 因为 \(1007=4\times252-1\)，代入 \(a_{4n-1}=0\) 得 \(a_{1007}=0\)，所以 \(a_{2014}=0\).
\end{solution}
\end{problem}

\section{解答题}

\begin{problem}
在 \(\triangle ABC\) 中，角 \(A\)，\(B\)，\(C\) 的对边分别为 \(a\)，\(b\)，\(c\)，\(B = \frac{\pi}{3}\)，\(\cos A = \frac{4}{5}\)，\(b = \sqrt{3}\).

\begin{enumerate}
\item 求 \(\sin C\) 的值；
\item 求 \(\triangle ABC\) 的面积.
\end{enumerate}
\end{problem}

\begin{solution}
\begin{enumerate}
\item 由 \(0<A<\pi\) 且 \(\cos A=\frac45>0\)，得 \(A\) 为锐角，从而 \(\sin A=\frac35\). 因为 \(C=\pi-(A+B)\)，所以 \(\sin C=\sin(A+B)=\sin A\cos B+\cos A\sin B=\frac35\cdot\frac12+\frac45\cdot\frac{\sqrt3}{2}=\frac{3+4\sqrt3}{10}\).
\item 由正弦定理，\(\dfrac{a}{\sin A}=\dfrac{b}{\sin B}=\dfrac{\sqrt3}{\sqrt3/2}=2\)，所以 \(a=2\sin A=\frac65\). 因而 \(\triangle ABC\) 的面积为 \(S=\frac12ab\sin C=\frac12\cdot\frac65\cdot\sqrt3\cdot\frac{3+4\sqrt3}{10}=\frac{36+9\sqrt3}{50}\).
\end{enumerate}
\end{solution}

\begin{problem}
如图，在三棱锥 \(P-ABC\) 中，\(PA \perp\) 底面 \(ABC\)，\(PA = AB\)，\(\angle ABC = 60^\circ\)，\(\angle BCA = 90^\circ\)，点 \(D\)，\(E\) 分别在棱 \(PB\)，\(PC\) 上，且 \(DE \parallel BC\).

\begin{enumerate}
\item 求证：\(BC \perp\) 平面 \(PAC\)；
\item 当 \(D\) 为 \(PB\) 的中点时，求 \(AD\) 与平面 \(PAC\) 所成的角的大小；
\item 是否存在点 \(E\) 使得二面角 \(A - DE - P\) 为直二面角？ 并说明理由.
\end{enumerate}

\texfigure[width=6.0cm]{img_repaint/2009/beijing_science/q16_fig1.tex}
\end{problem}

\begin{solution}
\begin{enumerate}
\item 因为 \(PA\perp\) 底面 \(ABC\)，所以 \(PA\perp BC\). 又因为 \(\angle BCA=90^{\circ}\)，所以 \(AC\perp BC\). 直线 \(PA\)，\(AC\) 是平面 \(PAC\) 内的两条相交直线，因此 \(BC\perp\) 平面 \(PAC\).
\item 令 \(PA=AB=1\). 在直角三角形 \(ABC\) 中，\(\angle B=60^{\circ}\)，\(\angle C=90^{\circ}\)，故 \(BC=\frac12\)，\(AC=\frac{\sqrt3}{2}\). 建立空间直角坐标系，取 \(C\) 为原点，令 \(BC\)，\(CA\)，\(AP\) 分别沿三个坐标轴，则可设 \(A\left(0,\frac{\sqrt3}{2},0\right)\)，\(B\left(\frac12,0,0\right)\)，\(P\left(0,\frac{\sqrt3}{2},1\right)\). 因为 \(D\) 为 \(PB\) 的中点，所以 \(D\left(\frac14,\frac{\sqrt3}{4},\frac12\right)\)，从而 \(\overrightarrow{AD}=\left(\frac14,-\frac{\sqrt3}{4},\frac12\right)\)，且 \(|AD|=\sqrt{\frac1{16}+\frac3{16}+\frac14}=\frac{\sqrt2}{2}\). 平面 \(PAC\) 的方程为 \(x=0\)，设 \(AD\) 与该平面所成的角为 \(\theta\)，则 \(AD\) 在平面法向上的分量为 \(\frac14\)，所以 \(\sin\theta=\frac{1/4}{\sqrt2/2}=\frac{\sqrt2}{4}\)，即 \(\theta=\arcsin\frac{\sqrt2}{4}=\arctan\frac1{\sqrt7}\).
\item 仍令 \(PA=AB=1\)，并设 \(t=\dfrac{PD}{PB}=\dfrac{PE}{PC}\). 由 \(DE\parallel BC\) 及三角形中平行线截线段成比例，确有上述两个比值相等. 于是 \(D\left(\frac{t}{2},\frac{\sqrt3(1-t)}{2},1-t\right)\)，\(E\left(0,\frac{\sqrt3(1-t)}{2},1-t\right)\)，其中 \(0<t<1\). 进一步有 \(\overrightarrow{ED}=\left(\frac{t}{2},0,0\right)\)，\(\overrightarrow{EA}=\left(0,\frac{\sqrt3t}{2},-(1-t)\right)\)，\(\overrightarrow{EP}=\left(0,\frac{\sqrt3t}{2},t\right)\). 直接计算向量积，得 \(\overrightarrow{ED}\times\overrightarrow{EA}=\left(0,\frac{t(1-t)}{2},\frac{\sqrt3t^2}{4}\right)\). 同理，\(\overrightarrow{ED}\times\overrightarrow{EP}=\left(0,-\frac{t^2}{2},\frac{\sqrt3t^2}{4}\right)\). 因为 \(0<t<1\)，可将两个向量同乘非零因子 \(\frac2t\)，得到平面 \(ADE\) 的一个法向量 \(\bs{n}_1=\left(0,1-t,\frac{\sqrt3t}{2}\right)\)，平面 \(PDE\) 的一个法向量为 \(\bs{n}_2=\left(0,-t,\frac{\sqrt3t}{2}\right)\).
两平面垂直当且仅当 \(\bs{n}_1\cdot\bs{n}_2=0\)，而 \(\bs{n}_1\cdot\bs{n}_2=-t(1-t)+\frac{3t^2}{4}=t\left(\frac{7t}{4}-1\right)\). 取 \(t=\frac47\) 即可使内积为零，此时 \(D\)，\(E\) 均在相应棱的内部，故存在点 \(E\) 使二面角 \(A-DE-P\) 为直二面角，且 \(\dfrac{PD}{PB}=\dfrac{PE}{PC}=\frac47\).
\end{enumerate}
\end{solution}

\begin{problem}
某学生在上学路上要经过 \(4\) 个路口，假设在各路口是否遇到红灯是相互独立的，遇到红灯的概率都是 \(\frac{1}{3}\)，遇到红灯时停留的时间都是 \(2 \mathrm{~min}\).

\begin{enumerate}
\item 求这名学生在上学路上到第三个路口时首次遇到红灯的概率；
\item 求这名学生在上学路上因遇到红灯停留的总时间 \(\xi\) 的分布列及期望.
\end{enumerate}
\end{problem}

\begin{solution}
\begin{enumerate}
\item 到第三个路口时首次遇到红灯，表示前两个路口未遇到红灯且第三个路口遇到红灯，所以所求概率为 \(\left(\frac23\right)^2\cdot\frac13=\frac4{27}\).
\item 设四个路口中遇到红灯的个数为 \(K\)，则 \(K\) 服从参数为 \(4\)，\(\frac13\) 的二项分布，且 \(\xi=2K\). 因此 \(\xi\) 的分布列为

\begin{center}
\begin{tabular}{c|ccccc}
\(\xi\) & \(0\) & \(2\) & \(4\) & \(6\) & \(8\) \\
\hline
\(P\) & \(\dfrac{16}{81}\) & \(\dfrac{32}{81}\) & \(\dfrac{8}{27}\) & \(\dfrac{8}{81}\) & \(\dfrac{1}{81}\)
\end{tabular}
\end{center}

其中，例如 \(P(\xi=2k)=\mathrm C_4^k\left(\frac13\right)^k\left(\frac23\right)^{4-k}\)，\(k=0\)，\(1\)，\(2\)，\(3\)，\(4\). 又 \(E(K)=4\cdot\frac13=\frac43\)，故 \(E(\xi)=2E(K)=\frac83\).
\end{enumerate}
\end{solution}

\begin{problem}
设函数 \( f(x) = x\e^{kx} \) （\(k \neq 0\)）.

\begin{enumerate}
\item 求曲线 \( y = f(x) \) 在点 \( (0, f(0)) \) 处的切线方程；
\item 求函数 \( f(x) \) 的单调区间；
\item 若函数 \( f(x) \) 在区间 \( (-1,1) \) 内单调递增，求 \(k\) 的取值范围.
\end{enumerate}
\end{problem}

\begin{solution}
\begin{enumerate}
\item \(f(0)=0\)，且 \(f'(x)=\e^{kx}(1+kx)\)，所以 \(f'(0)=1\). 因而曲线在 \((0,f(0))=(0,0)\) 处的切线方程为 \(y=x\).
\item 因为 \(\e^{kx}>0\)，所以 \(f'(x)\) 的符号由 \(1+kx\) 决定. 当 \(k>0\) 时，\(f'(x)<0\) 在 \(\left(-\infty,-\frac1k\right)\) 上成立，\(f'(x)>0\) 在 \(\left(-\frac1k,+\infty\right)\) 上成立，故函数在 \(\left(-\infty,-\frac1k\right)\) 上单调递减，在 \(\left(-\frac1k,+\infty\right)\) 上单调递增. 当 \(k<0\) 时，不等式方向相反，故函数在 \(\left(-\infty,-\frac1k\right)\) 上单调递增，在 \(\left(-\frac1k,+\infty\right)\) 上单调递减.
\item 要使函数在 \((-1,1)\) 内单调递增，需使 \(1+kx\geqslant0\) 对任意 \(x\in(-1,1)\) 成立. 当 \(k>0\) 时，这等价于 \(1-k\geqslant0\)，即 \(0<k\leqslant1\)；当 \(k<0\) 时，这等价于 \(1+k\geqslant0\)，即 \(-1\leqslant k<0\). 结合 \(k\neq0\)，得 \(k\in[-1,0)\cup(0,1]\).
\end{enumerate}
\end{solution}

\begin{problem}
已知双曲线 \(C: \frac{x^2}{a^2} - \frac{y^2}{b^2} = 1\) （\(a > 0\)，\(b > 0\)）的离心率为 \(\sqrt{3}\)，右准线方程为 \(x = \frac{\sqrt{3}}{3}\).

\begin{enumerate}
\item 求双曲线 \(C\) 的方程；
\item 设直线 \(l\) 是圆 \(O: x^{2} + y^{2} = 2\) 上动点 \(P(x_{0}, y_{0})\)（\(x_{0} y_{0} \neq 0\)）处的切线，\(l\) 与双曲线 \(C\) 交于不同的两点 \(A\)，\(B\)，证明 \(\angle AOB\) 的大小为定值.
\end{enumerate}
\end{problem}

\begin{solution}
\begin{enumerate}
\item 设双曲线的半焦距为 \(c\)，则 \(c^2=a^2+b^2\)，且离心率 \(e=\frac ca=\sqrt3\). 双曲线右准线方程为 \(x=\frac ae\)，所以 \(\frac{a}{\sqrt3}=\frac{\sqrt3}{3}\)，得 \(a=1\). 从而 \(c=\sqrt3\)，\(b^2=c^2-a^2=2\)，故双曲线 \(C\) 的方程为 \(x^2-\frac{y^2}{2}=1\).
\item 由 \(P(x_0,y_0)\) 在圆 \(x^2+y^2=2\) 上，且 \(y_0\neq0\)，圆在 \(P\) 处的切线可写为 \(x_0x+y_0y=2\)，进一步写成 \(l:y=rx+m\)，其中 \(r=-\frac{x_0}{y_0}\)，\(m=\frac2{y_0}\). 由 \(x_0^2+y_0^2=2\) 得 \(y_0^2(1+r^2)=2\)，所以 \(m^2=\frac4{y_0^2}=2(1+r^2)\). 设 \(A(x_1,y_1)\)，\(B(x_2,y_2)\)，则 \(y_i=rx_i+m\). 将直线方程代入双曲线方程，得 \((2-r^2)x^2-2rmx-(m^2+2)=0\). 因为交点 \(A\)，\(B\) 不同，二次项系数不为零，故由韦达定理，\(x_1+x_2=\frac{2rm}{2-r^2}\)，\(x_1x_2=-\frac{m^2+2}{2-r^2}\). 于是
\[
\overrightarrow{OA}\cdot\overrightarrow{OB}=x_1x_2+y_1y_2=(1+r^2)x_1x_2+rm(x_1+x_2)+m^2
=\frac{-(1+r^2)(m^2+2)+2r^2m^2}{2-r^2}+m^2
\]
由 \(m^2=2(1+r^2)\)，上式等于：
\[
\frac{2(1+r^2)(r^2-2)}{2-r^2}+2(1+r^2)=0
\]
因此 \(OA\perp OB\)，从而 \(\angle AOB=90^{\circ}\)，其大小为定值.
\end{enumerate}
\end{solution}

\begin{problem}
已知数集 \(A = \{a_{1}, a_{2}, \cdots, a_{n}\}\) （\(1 \leqslant a_{1} < a_{2} < \cdots < a_{n}\)，\(n \geqslant 2\)）具有性质 \(P\)：对任意的 \(i\)，\(j\) （\(1 \leqslant i \leqslant j \leqslant n\)），\(a_{i} a_{j}\) 与 \(\frac{a_{j}}{a_{i}}\) 两数中至少有一个属于 \(A\).

\begin{enumerate}
\item 分别判断集数 \(\{1, 3, 4\}\) 与 \(\{1, 2, 3, 6\}\) 是否具有性质 \(P\)，并说明理由；
\item 证明：\(a_{1}=1\)，且 \(\frac{a_{1}+a_{2}+\cdots+a_{n}}{a_{1}^{-1}+a_{2}^{-1}+\cdots+a_{n}^{-1}}=a_{n}\)；
\item 证明：当 \(n = 5\) 时，\(a_1\)，\(a_2\)，\(a_3\)，\(a_4\)，\(a_5\) 成等比数列.
\end{enumerate}
\end{problem}

\begin{solution}
\begin{enumerate}
\item 对集合 \(\{1,3,4\}\)，取 \(a_i=3\)，\(a_j=4\)，则 \(a_i a_j=12\notin\{1,3,4\}\)，且 \(\frac{a_j}{a_i}=\frac43\notin\{1,3,4\}\)，所以该集合不具有性质 \(P\). 对集合 \(\{1,2,3,6\}\)，含有元素 \(1\) 的各对中，乘积仍属于该集合；其余各对分别满足 \(\frac22=1\)，\(2\cdot3=6\)，\(\frac62=3\)，\(\frac33=1\)，\(\frac63=2\)，\(\frac66=1\)，所以该集合具有性质 \(P\).
\item 先证 \(a_1=1\). 若 \(a_1>1\)，则 \(1\notin A\). 对 \((a_1,a_1)\) 应用性质 \(P\)，只能有 \(a_1^2\in A\)；再对 \((a_1^2,a_1^2)\) 应用同样的论证，得到 \(a_1^4\in A\)，如此继续可得 \(a_1^{2^m}\in A\) 对任意正整数 \(m\) 成立，这与 \(A\) 是有限集矛盾，故 \(a_1=1\). 对任意 \(i\geqslant2\)，将 \((a_i,a_n)\) 代入性质 \(P\). 因为 \(a_i>1\)，所以 \(a_i a_n>a_n\)，不可能属于 \(A\)，从而 \(\frac{a_n}{a_i}\in A\). 对 \(i=1\)，也有 \(\frac{a_n}{a_1}=a_n\in A\). 因此 \(\left\{\frac{a_n}{a_i}:1\leqslant i\leqslant n\right\}=A\)，于是 \(a_n\sum_{i=1}^n\frac1{a_i}=\sum_{i=1}^n\frac{a_n}{a_i}=\sum_{i=1}^n a_i\). 分母为正，故 \(\frac{a_1+a_2+\cdots+a_n}{a_1^{-1}+a_2^{-1}+\cdots+a_n^{-1}}=a_n\).
\item 令 \(n=5\)，并记 \(x=a_2\)，\(y=a_3\)，\(z=a_4\)，\(w=a_5\). 由第（2）问，\(\left\{w,\frac{w}{x},\frac{w}{y},\frac{w}{z},1\right\}=A\). 因为 \(1<x<y<z<w\)，左侧元素按从大到小排列为 \(w\)，\(\frac{w}{x}\)，\(\frac{w}{y}\)，\(\frac{w}{z}\)，\(1\)，而 \(A\) 按从小到大排列为 \(1\)，\(x\)，\(y\)，\(z\)，\(w\)，所以 \(\frac{w}{x}=z\)，\(\frac{w}{y}=y\)，即 \(xz=w=y^2\). 因为 \(y<z<w=y^2\)，所以 \(yz>w\)，对 \((a_3,a_4)\) 应用性质 \(P\) 时，乘积不在 \(A\)，只能有 \(\frac{z}{y}\in A\). 又 \(1<\frac{z}{y}<y\)，而 \(A\) 中严格位于 \(1\) 与 \(y=a_3\) 之间的唯一元素是 \(x=a_2\)，故 \(\frac{z}{y}=x\). 结合 \(xz=y^2\)，得 \(x^2=y\)，进而 \(z=x^3\)，\(w=x^4\). 因此 \(a_1=1\)，\(a_2=x\)，\(a_3=x^2\)，\(a_4=x^3\)，\(a_5=x^4\)，成等比数列.
\end{enumerate}
\end{solution}
